\documentclass[12pt]{article}

\usepackage{sbc-template}
\usepackage{graphicx,url}
\usepackage[utf8]{inputenc}
\usepackage[brazil]{babel}
\usepackage{booktabs}
\usepackage{hyperref}

\sloppy

\title{Front-end Web Integrado à API RESTful e Banco de Dados Relacional para Gestão de Locadora de Veículos}

\author{Gabriel [Sobrenome]\inst{1}}

\address{Curso de Análise e Desenvolvimento de Sistemas -- Tecnologias para Análise e Desenvolvimento de Sistemas\\
  [Nome da Instituição] -- [Cidade/UF] -- Brasil
  \email{seu.email@instituicao.edu.br}
}

\begin{document}

\maketitle

\begin{abstract}
  This paper describes the prototyping and development of a web front-end application integrated with the RESTful API and relational database created in the first part of the Vehicle Rental System (\textit{LocadoraVeiculos}). Built with React, TypeScript, Vite, and Tailwind CSS, the application consumes all CRUD endpoints and analytical routes based on \texttt{INNER JOIN} and \texttt{LEFT OUTER JOIN} queries executed in Microsoft SQL Server via ASP.NET Core (.NET 9) and Entity Framework Core. All screens operate with live database records, supporting multi-field filtering for each entity and visual reports.
\end{abstract}

\begin{resumo}
  Este artigo apresenta a prototipação e o desenvolvimento de uma aplicação front-end web integrada à API RESTful e ao banco de dados relacional construídos na primeira parte do Sistema de Aluguel de Veículos (\textit{LocadoraVeiculos}). Desenvolvido com React, TypeScript, Vite e Tailwind CSS, o front-end consome todas as rotas CRUD e as consultas com \texttt{INNER JOIN} e \texttt{LEFT OUTER JOIN} executadas no Microsoft SQL Server por meio do ASP.NET Core (.NET 9) e Entity Framework Core. Todas as telas operam com dados reais do banco, oferecendo filtros em múltiplos campos por entidade e relatórios visuais.
\end{resumo}

\section{Introdução}

O controle de locação de veículos envolve o acompanhamento simultâneo da disponibilidade da frota, do cadastro de clientes, da vigência de contratos e do cálculo de valores na devolução. Quando realizado sem integração entre interface e banco de dados, o processo fica sujeito a reservas duplicadas, inconsistências de quilometragem e dificuldade na análise de faturamento.

Na primeira parte deste projeto (Trabalho Prático 1), foi desenvolvido o back-end \texttt{LocadoraVeiculos.API} em C\# (.NET 9), usando ASP.NET Core Web API, Entity Framework Core (Code-First) e Microsoft SQL Server. A API disponibilizou operações CRUD para cinco entidades e cinco rotas de consulta usando junções relacionais (\texttt{INNER JOIN} e \texttt{LEFT OUTER JOIN}).

Nesta segunda etapa (Trabalho Prático 2), foi desenvolvido o front-end \texttt{LocadoraVeiculos.Web}, conectado diretamente ao back-end anterior. O trabalho abrange a criação de wireframes, a implementação de telas de cadastro, edição, exclusão e pesquisa com pelo menos dois filtros por entidade, além da visualização gráfica e tabular das consultas relacionais.

\section{Modelagem e Arquitetura}

O sistema segue o modelo cliente-servidor com comunicação HTTP/JSON:

\begin{itemize}
  \item \textbf{Banco de Dados (SQL Server / LocalDB):} formado pelas tabelas \texttt{Fabricantes}, \texttt{Categorias}, \texttt{Veiculos}, \texttt{Clientes} e \texttt{Alugueis}, com chaves estrangeiras configuradas com \texttt{OnDelete: Restrict} e restrições de unicidade para nome de fabricante, nome de categoria, placa, CPF e e-mail.
  \item \textbf{Back-end (\texttt{LocadoraVeiculos.API}):} API em ASP.NET Core (.NET 9) organizada em \textit{Models}, \textit{Data}, \textit{DTOs}, \textit{Controllers} e \textit{ExceptionMiddleware}. Para integração com o navegador e suporte aos filtros de pesquisa, foi configurada política de CORS em \texttt{Program.cs} e foram adicionados parâmetros opcionais (\texttt{[FromQuery]}) nos métodos \texttt{GET} dos controladores CRUD, mantendo compatibilidade com as rotas da etapa anterior.
  \item \textbf{Front-end (\texttt{LocadoraVeiculos.Web}):} aplicação em React 18, TypeScript, Vite e Tailwind CSS, utilizando \texttt{react-router-dom} para navegação, \texttt{Recharts} para gráficos e um módulo centralizado (\texttt{src/services/api.ts}) para requisições HTTP e tratamento de respostas de erro (\texttt{400}, \texttt{404}, \texttt{409} e \texttt{500}).
\end{itemize}

\section{Wireframes}

Antes da implementação das telas, foram elaborados wireframes de média/alta fidelidade definindo o menu lateral fixo, os cards de indicadores, as barras de filtro, as tabelas de listagem e os modais de cadastro, edição, histórico e devolução, cobrindo as sete telas principais: \textit{Dashboard}, \textit{Fabricantes}, \textit{Categorias}, \textit{Veículos}, \textit{Clientes}, \textit{Aluguéis} e \textit{Consultas}.

\begin{itemize}
  \item \textbf{Link dos wireframes no Figma:} \url{https://www.figma.com/file/INSIRA-SEU-LINK-AQUI}
  \item \textbf{Referência local no repositório:} \texttt{Docs/WIREFRAMES\_GUIA.md} e \texttt{Docs/wireframe\_prototipo.html}.
\end{itemize}

\section{Implementação e Integração Front-end e API}

Cada módulo da aplicação comunica-se diretamente com as rotas criadas na API, sem dados fixos no cliente. A Tabela~\ref{tab:rotas} resume as rotas utilizadas e os filtros disponíveis por entidade.

\begin{table}[ht]
\centering
\caption{Relação entre módulos do front-end, rotas da API e filtros implementados.}
\label{tab:rotas}
\scriptsize
\begin{tabular}{@{}p{2.3cm}p{5.5cm}p{5.5cm}@{}}
\toprule
\textbf{Entidade} & \textbf{Rotas Consumidas} & \textbf{Filtros Personalizados} \\
\midrule
Fabricantes & \texttt{GET/POST/PUT/DELETE /api/fabricantes}, \texttt{GET /api/consultas/fabricantes-desempenho} & Nome; País de origem. \\
Categorias & \texttt{GET/POST/PUT/DELETE /api/categorias}, \texttt{GET /api/consultas/faturamento-por-categoria} & Nome/descrição; Diária mínima; Diária máxima. \\
Veículos & \texttt{GET/POST/PUT/DELETE /api/veiculos}, \texttt{GET /api/veiculos/disponiveis}, \texttt{GET /api/consultas/veiculos-utilizacao} & Modelo; Placa; Fabricante; Categoria; Disponibilidade. \\
Clientes & \texttt{GET/POST/PUT/DELETE /api/clientes}, \texttt{GET /api/alugueis/cliente/\{id\}}, \texttt{GET /api/consultas/clientes-resumo-locacoes} & Nome; CPF; E-mail; Situação de locação. \\
Aluguéis & \texttt{GET/POST /api/alugueis}, \texttt{PUT /api/alugueis/\{id\}/devolver}, \texttt{DELETE /api/alugueis/\{id\}}, \texttt{GET /api/consultas/alugueis-detalhados} & Cliente; Modelo; Fabricante; Status; Data inicial e final. \\
\bottomrule
\end{tabular}
\end{table}

\subsection{Fluxo de Locação e Validações}

Na abertura de uma nova locação, o formulário consulta \texttt{GET /api/veiculos/disponiveis} para exibir apenas veículos livres. Ao confirmar (\texttt{POST /api/alugueis}), o back-end altera o status do veículo para indisponível. Na devolução (\texttt{PUT /api/alugueis/\{id\}/devolver}), o servidor valida a data e a quilometragem final, calcula o valor total conforme as diárias, atualiza a quilometragem do veículo e libera o automóvel para novos contratos. Tentativas de cadastro duplicado (\texttt{409}) ou exclusão de registros com vínculos (\texttt{400}) retornam a mensagem da API diretamente em alertas na tela.

\subsection{Consultas com INNER JOIN e LEFT OUTER JOIN}

A página de Consultas reúne as cinco rotas analíticas do \texttt{ConsultasController}:
\begin{enumerate}
  \item \textbf{Aluguéis Detalhados:} \texttt{INNER JOIN} entre \texttt{Alugueis}, \texttt{Clientes}, \texttt{Veiculos}, \texttt{Fabricantes} e \texttt{Categorias}.
  \item \textbf{Faturamento por Categoria:} \texttt{INNER JOIN} com agrupamento e soma para cálculo de faturamento total e ticket médio, exibido em tabela e gráfico de barras.
  \item \textbf{Resumo de Clientes:} \texttt{LEFT OUTER JOIN} entre \texttt{Clientes} e \texttt{Alugueis}, permitindo filtrar inclusive clientes sem nenhuma locação registrada.
  \item \textbf{Utilização de Veículos:} junção mista (\texttt{INNER JOIN} e \texttt{LEFT OUTER JOIN}) com contagem de locações, quilometragem rodada e receita por veículo.
  \item \textbf{Desempenho por Fabricante:} múltiplo \texttt{LEFT OUTER JOIN} entre \texttt{Fabricantes}, \texttt{Veiculos} e \texttt{Alugueis}, apresentado em gráfico e tabela.
\end{enumerate}

\section{Resultados}

Os testes foram realizados com o back-end (.NET 9) e o front-end (Vite/React) em execução local sobre o banco SQL Server LocalDB. Foi validado o ciclo completo de cadastro, edição, exclusão e filtragem das cinco entidades, o tratamento de erros de integridade e unicidade, bem como a atualização imediata dos gráficos do painel e das tabelas de consultas após o registro e a devolução de contratos.

\section{Conclusão}

Com o desenvolvimento do front-end \texttt{LocadoraVeiculos.Web}, o sistema passou a oferecer interface completa para operação e consulta dos dados modelados na Parte 1. A definição prévia dos wireframes facilitou a construção das telas e a separação entre camada de serviço HTTP e componentes garantiu aderência aos contratos da API. O código-fonte completo e o histórico de commits estão disponíveis em: \url{https://github.com/SEU-USUARIO/SEU-REPOSITORIO}.

\bibliographystyle{sbc}

\end{document}
